\documentclass[11pt,a4paper]{article}

% --- Paquetes de Formato y Estilo ---
\usepackage[utf8]{inputenc}
\usepackage[spanish]{babel}
\usepackage[margin=2.5cm]{geometry}
\usepackage{graphicx}
\usepackage{booktabs}
\usepackage{amsmath}
\usepackage{amsfonts}
\usepackage{xcolor}
\usepackage{hyperref}
\usepackage{enumitem}

% --- Configuración de Enlaces ---
\hypersetup{
    colorlinks=true,
    linkcolor=blue,
    filecolor=magenta,      
    urlcolor=cyan,
    pdftitle={Estructura y Contenido de las Presentaciones - Proyecto Final},
    pdfpagemode=FullScreen,
}

% --- Colores Premium ---
\definecolor{UtecCyan}{RGB}{0, 191, 255}       % Azul cian
\definecolor{UtecDarkGray}{RGB}{50, 50, 50}    % Gris Oscuro
\definecolor{LightGray}{RGB}{245, 245, 245}

% --- Metadatos del Documento ---
\title{\textbf{Estructura y Contenido de las Presentaciones Finales}\\
\large Maestría en Ciencia de Datos e Inteligencia Artificial -- UTEC Posgrado}
\author{\textbf{Grupo G2 (Sección 1 - g102)}\\ 
Integrantes: David Jimenez, Herles Pinedo, Jose Nomberto, Kevin Luyo, Yessenia Chirinos}
\date{18 de Julio de 2026}

\begin{document}

\maketitle
\tableofcontents
\newpage

% -------------------------------------------------------------------
\section{Introducción y Propósito}
% -------------------------------------------------------------------
Este documento técnico detalla la estructura y el contenido exacto diapositiva por diapositiva para los entregables de presentación del proyecto final. Incorpora los componentes de automatización mediante \textbf{Databricks Jobs (Workflows)}, la gobernanza sobre \textbf{Unity Catalog} y la documentación detallada de la capa Gold para el grupo \textbf{g102}.

Siguiendo las directrices del proyecto final, toda mención de la mina real ha sido anonimizada bajo la denominación institucional de \textbf{Compañía Minera Los Andes S.A.} (o \textit{Minera Los Andes}).

---

% -------------------------------------------------------------------
\section{Parte A: Presentación Ejecutiva (Negocio y Storytelling)}
% -------------------------------------------------------------------
\textit{Enfoque: Directo, conceptual, visual, sin código. Alineado estrictamente a la estructura oficial de UTEC Posgrado: \textbf{Why -- What -- How -- When}.}

\subsection{Diapositiva 1: Carátula}
\begin{itemize}[label=$\bullet$]
    \item \textbf{Título:} Pipeline de Ingesta y Exposición de Data Products para Control de Inventario en Consignación.
    \item \textbf{Subtítulo:} Proyecto Integrador de Ingeniería de Datos (Etapa Final).
    \item \textbf{Institución:} Universidad de Ingeniería y Tecnología -- UTEC Posgrado.
    \item \textbf{Grupo:} Grupo G2 (Sección 1 - g102).
    \item \textbf{Integrantes:} David Jimenez, Herles Pinedo, Jose Nomberto, Kevin Luyo, Yessenia Chirinos.
\end{itemize}

\subsection{Diapositiva 2: El Contexto (El Negocio de Minera Los Andes)}
\begin{itemize}[label=$\bullet$]
    \item \textbf{Título:} Logística de Repuestos Críticos y Capital de Trabajo.
    \item \textbf{Mensaje Clave:} En operaciones mineras a tajo abierto a gran escala, la disponibilidad de herramientas, consumibles de seguridad (EPP) y repuestos mecánicos es vital para asegurar la continuidad operativa de la planta.
    \item \textbf{Consignación:} Los materiales se almacenan físicamente en la mina, pero pertenecen legalmente a los proveedores (Ferreyros, Komatsu, Shell, Epiroc, Moly-Cop) hasta que la mina los consume, momento en el cual se efectúa el cobro y reposición.
\end{itemize}

\subsection{Diapositiva 3: WHY (¿Por qué estamos haciendo esto?)}
\begin{itemize}[label=$\bullet$]
    \item \textbf{Título:} El Desafío Financiero y Operativo.
    \item \textbf{Dolores de Negocio:}
    \begin{itemize}[label=--]
        \item \textbf{Stockouts (Desabastecimiento):} La falta de un repuesto crítico paraliza la planta o frentes de minado, con pérdidas estimadas de \textbf{\$10,000 a \$50,000 USD por hora}.
        \item \textbf{Falta de visibilidad:} El control de inventarios se realiza en SAP local. Los proveedores reciben reportes manuales por correo de forma semanal/quincenal, impidiendo una reacción proactiva.
        \item \textbf{Consecuencias:} Sobrestock preventivo ineficiente en los almacenes físicos de la mina o paradas de planta imprevistas.
    \end{itemize}
\end{itemize}

\subsection{Diapositiva 4: HOW (¿Cómo lo estamos solucionando?)}
\begin{itemize}[label=$\bullet$]
    \item \textbf{Título:} Arquitectura de Datos Desacoplada y Distribuida.
    \item \textbf{Arquitectura de Ingesta (Airflow + Azure):} Ingesta automática diaria e inmutable de los reportes SAP en local a Azure Data Lake Storage (capa Raw, ruta: \texttt{raw/airflow/G2}).
    \item \textbf{Orquestación en Databricks Workflows:} Automatización end-to-end estructurada en tareas secuenciales programadas diariamente (Bronze $\rightarrow$ Silver $\rightarrow$ Gold $\rightarrow$ Refresco/Auditoría).
    \item \textbf{Arquitectura Medallion en Databricks:}
    \begin{itemize}[label=--]
        \item \textbf{Bronze:} Consolidación histórica 1:1 de los archivos.
        \item \textbf{Silver:} Calidad (deduplicación, enmascaramiento PII) y enriquecimiento con el maestro de materiales; desvío automático de errores a \textbf{Cuarentena}.
        \item \textbf{Gold:} Modelado de negocio y consolidación de los \textbf{3 Data Products}.
    \end{itemize}
\end{itemize}

\subsection{Diapositiva 5: WHAT (¿Qué construimos?)}
\begin{itemize}[label=$\bullet$]
    \item \textbf{Título:} Descentralización a través de Data Products.
    \item \textbf{Data Product 1 - Alertas de Reabastecimiento Crítico (\texttt{gold\_replenishment\_alerts}):} Calcula el stock actual agrupando ingresos y consumos de consignación. Dispara alertas de compra si baja del punto de reorden.
    \item \textbf{Data Product 2 - Análisis de Consumo y Costos (\texttt{gold\_consumption\_analysis}):} Visualiza consumos mensuales históricos para renegociar contratos.
    \item \textbf{Data Product 3 - Conciliación Mensual (\texttt{gold\_monthly\_conciliation}):} Detalla los consumos facturables de fin de mes, automatizando el proceso de conciliación contable.
\end{itemize}

\subsection{Diapositiva 6: WHAT (Aplicación de Data Mesh y Gobierno del Dato)}
\begin{itemize}[label=$\bullet$]
    \item \textbf{Título:} Gobierno del Dato y Descentralización de la Propiedad.
    \item \textbf{Propiedad por Dominio:} Catálogos independientes creados en Unity Catalog mediante el registro oficial de la clase: \texttt{g102\_log\_reabastecimiento}, \texttt{g102\_log\_consumos}, \texttt{g102\_fin\_conciliacion}.
    \item \textbf{Gobernanza Computacional Federada:}
    \begin{itemize}[label=--]
        \item \textbf{Seguridad PII:} Enmascaramiento dinámico a nivel de columna (UDF SQL) sobre el campo \texttt{User Name} (Usuario SAP de mina) para restringir el acceso a no administradores.
        \item \textbf{Contratos de Datos:} Definición automática de esquemas y reglas de calidad en YAML para garantizar la fiabilidad del Data Product.
    \end{itemize}
\end{itemize}

\subsection{Diapositiva 7: WHAT (Valor de Negocio Generado)}
\begin{itemize}[label=$\bullet$]
    \item \textbf{Título:} Impacto en la Cadena de Abastecimiento.
    \item \textbf{Reducción de Paradas de Planta:} Alerta proactiva que reduce el riesgo de stockouts críticos en un \textbf{90\%}.
    \item \textbf{Eficiencia Financiera:} Reducción del sobrestock preventivo en almacén, liberando espacio físico y optimizando el capital de trabajo de los proveedores.
    \item \textbf{Productividad Administrativa:} Ahorro de más de \textbf{20 horas semanales} de trabajo manual de conciliación por parte del equipo logístico y financiero.
\end{itemize}

\subsection{Diapositiva 8: WHEN (Estado Actual y Hoja de Ruta)}
\begin{itemize}[label=$\bullet$]
    \item \textbf{Título:} Plan de Despliegue y Resultados del MVP.
    \item \textbf{Hitos Completados:}
    \begin{itemize}[label=--]
        \item Ingesta automática local a nube (Airflow DAG estable en Azure).
        \item Procesamiento Medallion y generación de los 3 Data Products en Databricks mediante vistas en Gold.
        \item Publicación de tablas y vistas en Unity Catalog.
        \item Monitoreo de SLAs simulado en base a consultas de uso.
    \end{itemize}
    \item \textbf{Hitos por Ejecutar (Q3 2026):}
    \begin{itemize}[label=--]
        \item Conexión directa del catálogo de Databricks con Power BI para tableros dinámicos en mina.
        \item Capacitación a los planificadores de proveedores en el consumo de las vistas.
    \end{itemize}
\end{itemize}

\subsection{Diapositiva 9: Lecciones Aprendidas y Desafíos}
\begin{itemize}[label=$\bullet$]
    \item \textbf{Título:} Retos Superados en el Camino.
    \item \textbf{Calidad de Origen:} SAP presentaba registros con delimitadores inconsistentes y nulos. Se resolvió con reglas de cuarentena y try-casts resilientes.
    \item \textbf{Gobernanza:} Aislamiento y control de seguridad para que ningún proveedor pueda visualizar información de sus competidores.
    \item \textbf{Data Mesh:} La transición cultural de tratar los datos como un reporte estático a tratarlos como un \textbf{Producto de Datos} útil y gobernado.
\end{itemize}

\subsection{Diapositiva 10: Cierre}
\begin{itemize}[label=$\bullet$]
    \item \textbf{Texto:} ¡Muchas gracias! ¿Tienen alguna pregunta?
    \item \textbf{Contacto:} Equipo de Ingeniería de Datos -- Minera Los Andes.
\end{itemize}

\newpage
% -------------------------------------------------------------------
\section{Parte B: Estructura de la Presentación Técnica}
% -------------------------------------------------------------------
\textit{Enfoque: Detalles de implementación, código, diagramas de arquitectura y especificaciones del modelo.}

\subsection{Estructura Diapositiva por Diapositiva}
\begin{enumerate}
    \item \textbf{Slide 1: Carátula e Introducción Técnica:} Título oficial y alcance técnico (Arquitectura Medallion y Gobernanza en Databricks).
    \item \textbf{Slide 2: Arquitectura del Pipeline:} Diagrama completo (ADLS Gen2 $\rightarrow$ Bronze $\rightarrow$ Silver $\rightarrow$ Gold $\rightarrow$ Cuarentena).
    \item \textbf{Slide 3: Capa Bronze (Ingesta Inmutable):} Ingesta con llaves SAS en Azure y normalización de columnas Delta.
    \item \textbf{Slide 4: Capa Silver (Calidad y Resiliencia):} Reglas de calidad y try-casts resilientes para evitar fallos por datos corruptos de SAP.
    \item \textbf{Slide 5: Desafíos Técnicos de SAP y Soluciones:} Tabla unificada de problemas (cuarentena por ceros a la izquierda, stock negativo por saldo de apertura y formateo regional decimal).
    \item \textbf{Slide 6: Seguridad PII en Unity Catalog:} UDF de enmascaramiento dinámico de usuarios basada en roles (\texttt{is\_member}).
    \item \textbf{Slide 7: Capa Gold (Data Products):} Implementación SQL de las vistas de Reabastecimiento, Consumo y Conciliación Mensual de Pagos.
    \item \textbf{Slide 8: Gobernanza y Data Contracts:} Diccionario de datos en Unity Catalog y autogeneración del contrato YAML.
    \item \textbf{Slide 9: Reporte de SLAs y Auditoría:} Monitoreo de latencia y tasas de éxito simulados por restricciones de cuenta en Unity Catalog.
    \item \textbf{Slide 10: Orquestación con Databricks Workflows:} Job secuencial programado diariamente a las 01:00 AM.
    \item \textbf{Slide 11: Oportunidades de Mejora de Arquitectura:} Auto Loader, orquestación híbrida remota, multi-workspace Delta Sharing y OpenLineage.
    \item \textbf{Slide 12: Cierre:} Agradecimiento y ronda de preguntas.
\end{enumerate}

\end{document}
